% DERIVED from formal_layer.json — read-only, do not edit.
\begin{assumptionv}{ass:holder-derivative-bound}[Uniform bound on derivatives]
\[
\max_{|\alpha|\le m}\|D^\alpha g\|_\infty\le L.
\]
\end{assumptionv}

\begin{assumptionv}{ass:holder-modulus}[Hölder modulus of derivatives]
\[
\max_{|\alpha|=m}\sup_{x\ne x'}
\frac{|D^\alpha g(x)-D^\alpha g(x')|}
{\|x-x'\|_\infty^{\beta-m}}
\le L.
\]
\end{assumptionv}

\begin{definitionv}{def:holder}[Intrinsic cube Hölder ball]
The intrinsic cube Hölder ball is
\[
\mathcal H^\beta(L)
:=
\left\{
g:\mathcal X\to\mathbb R:
g\in C^m(\mathcal X),\
g\text{ satisfies }\cref{obj:ass:holder-derivative-bound,obj:ass:holder-modulus}
\right\},
\qquad
\mathcal X=[0,1]^d,
\qquad
m=\max\{\lceil\beta\rceil-1,0\}.
\]
Here \(d\) is a nonnegative integer and \(\beta,L\in\mathbb R\). The integer \(m\) is the derivative order, and the top-order modulus has exponent \(\beta-m\). The class \(C^m(\mathcal X)\) consists of functions whose classical coordinate partial derivatives through order \(m\) on the interior of the cube admit unique continuous traces on \(\mathcal X\). Every derivative \(D^\alpha g\), supremum, and difference quotient in the defining restrictions uses these traces. At positive integer \(\beta\), this convention uses Lipschitz derivatives of order \(\beta-1\).
\end{definitionv}

\begin{assumptionv}{ass:iid}[Independent and identically distributed observations]
The observed sample \(\mathcal D_n\) consists of \(n\) independent observations with common law \(P\), so that
\[
\mathcal D_n\sim P^{\otimes n}.
\]
\end{assumptionv}

\begin{assumptionv}{ass:density}[Covariate density lower bound]
The covariate density \(f\) satisfies
\[
c_f\le f(x)
\]
for Lebesgue-almost every \(x\in\mathcal X\).
\end{assumptionv}

\begin{assumptionv}{ass:consistency}[Consistency of observed outcomes]
The observed outcome \(Y\), binary treatment indicator \(A\), and potential outcomes \(Y(0)\) and \(Y(1)\) satisfy
\[
Y=AY(1)+(1-A)Y(0)
\]
almost surely.
\end{assumptionv}

\begin{assumptionv}{ass:exchangeability}[Conditional exchangeability of treatment]
The potential outcomes \(Y(0)\) and \(Y(1)\) are jointly independent of the treatment indicator \(A\) conditional on the covariate vector \(X\):
\[
(Y(0),Y(1))\perp A\mid X.
\]
\end{assumptionv}

\begin{assumptionv}{ass:global-tail}[Global propensity tail bound]
The propensity \(e(X)\), evaluated at the covariate vector \(X\), satisfies
\[
P\{e(X)\le t\}\le Ct^{\gamma-1}
\qquad\text{for every }t\in[0,1].
\]
\end{assumptionv}

\begin{assumptionv}{ass:bounded-potential-outcomes}[Treated response magnitude bound]
The observed outcome \(Y\), binary treatment indicator \(A\), and covariate vector \(X\) satisfy
\[
\operatorname*{ess\,sup}_{x\sim P_X}
\left(
 \left|\mathbb E_P[Y\mid A=1,X=x]\right|
 \vee
 \sup_{\lambda\in\mathbb R\setminus\{0\}}
 \frac{\sqrt{2\log \mathbb E_P\!\left[
   \exp\!\left\{\lambda\bigl(Y-\mathbb E_P[Y\mid A=1,X]\bigr)\right\}
   \mid A=1,X=x\right]}}{|\lambda|}
\right)
\le B.
\]
\end{assumptionv}

\begin{assumptionv}{ass:holder-response}[Hölder smoothness of response]
The treated causal response \(\mu_1\) belongs to the intrinsic cube Hölder ball \(\mathcal H^\beta(L)\):
\[
\mu_1\in\mathcal H^\beta(L).
\]
\end{assumptionv}

\begin{definitionv}{def:model}[Observational causal law class $\mathcal P_\gamma$]
The class \(\mathcal P_\gamma\) consists of observational causal laws \(P\) satisfying the following restrictions, with common fixed constants:
\begin{itemize}
\item \textbf{(Covariate density.)} The density \(f\) of the covariate vector \(X\) satisfies \cref{obj:ass:density}.
\item \textbf{(Consistency.)} The observed outcome \(Y\), binary treatment indicator \(A\), and potential outcomes \(Y(0)\) and \(Y(1)\) satisfy \cref{obj:ass:consistency}.
\item \textbf{(Exchangeability.)} The potential outcomes and treatment satisfy \cref{obj:ass:exchangeability}.
\item \textbf{(Overlap.)} The propensity \(e(X)\) satisfies \cref{obj:ass:global-tail}.
\item \textbf{(Response bound.)} The treated conditional outcome distribution satisfies \cref{obj:ass:bounded-potential-outcomes}.
\item \textbf{(Response smoothness.)} The treated causal response \(\mu_1\) satisfies \cref{obj:ass:holder-response}.
\end{itemize}
\end{definitionv}

\begin{definitionv}{def:risks}[Minimax risks $R_{\mathrm{pt}}$ and $R_\infty$]
The minimax expected absolute error at \(x_0\) and the minimax expected spatial-supremum error are, respectively,
\[
R_{\mathrm{pt}}(n,\gamma,x_0)
=
\inf_T\sup_{P\in\mathcal P_\gamma}
\mathbb E_P\left|T-\mu_{1,P}(x_0)\right|
\]
and
\[
R_\infty(n,\gamma)
=
\inf_T\sup_{P\in\mathcal P_\gamma}
\mathbb E_P\left[
\sup_{x\in\mathcal X}\left|T(x)-\mu_{1,P}(x)\right|
\right].
\]
Here \(\mu_{1,P}\) is the treated response under law \(P\), and each infimum ranges over estimators measurable with respect to the observed sample \(\mathcal D_n\). The first infimum uses scalar-valued estimators, and the second uses function-valued estimators.
\end{definitionv}

\begin{definitionv}{synth_1}[Tensor interpolation nodes]
The tensor node \(v_\ell\in\mathbb R^d\), for nonnegative integers \(d,m\) and an index vector \(\ell\in\{0,\ldots,m\}^d\), is defined by
\[
(v_\ell)_i=\frac{\ell_i+1}{m+2},
\qquad i=1,\ldots,d.
\]
\end{definitionv}

\begin{definitionv}{synth_2}[Polynomial template and microcell width]
Fix a dimension \(d\ge1\) and a polynomial degree \(m\in\mathbb N_0\). Let
\[
\mathcal I_m:=\{\alpha\in\mathbb N_0^d:|\alpha|\le m\},
\qquad p:=|\mathcal I_m|,
\qquad U(z):=(z^\alpha)_{\alpha\in\mathcal I_m},
\quad z\in[0,1]^d,
\]
with the coordinates of \(U\) ordered lexicographically. For
\(\ell\in\{0,\ldots,m\}^d\), define
\[
v_\ell:=\left(\frac{\ell_1+1}{m+2},\ldots,
\frac{\ell_d+1}{m+2}\right),
\qquad J:=(m+1)^d,
\]
and set
\[
G_0:=\frac1J\sum_{\ell\in\{0,\ldots,m\}^d}
U(v_\ell)U(v_\ell)^\top,
\qquad \lambda_0:=\lambda_{\min}(G_0).
\]
The template Lipschitz constant is
\[
L_U:=
\sup_{\substack{z,z'\in[0,1]^d\\z\ne z'}}
\frac{\|U(z)U(z)^\top-U(z')U(z')^\top\|_{\mathrm{op}}}
{\|z-z'\|_\infty},
\]
where \(\|\cdot\|_{\mathrm{op}}\) is the matrix operator norm induced by the Euclidean vector norm. This finite constant gives the bound
\[
\|U(z)U(z)^\top-U(z')U(z')^\top\|_{\mathrm{op}}
\le L_U\|z-z'\|_\infty.
\]
Define the microcell width and reference microcubes by
\[
\eta:=\min\left\{\frac1{2(m+2)},\frac{\lambda_0}{1+L_U}\right\},
\qquad
V_\ell:=v_\ell+[-\eta/2,\eta/2]^d.
\]
In particular, for \(z\in V_\ell\), the bound entering the width rule is
\[
\|U(z)U(z)^\top-U(v_\ell)U(v_\ell)^\top\|_{\mathrm{op}}
\le \frac{L_U\eta}{2}\le\frac{\lambda_0}{2}.
\]
\end{definitionv}

\begin{definitionv}{def:template}[Polynomial template $G_0$ and microcubes $V_\ell$]
The equally weighted template Gram matrix \(G_0\) and its least eigenvalue \(\lambda_0\) are
\[
G_0=\frac{1}{J}\sum_{\ell=1}^{J}U(v_\ell)U(v_\ell)^\top,
\qquad
\lambda_0=\lambda_{\min}(G_0),
\]
where \(U\) is the polynomial template vector, \(v_\ell\) are its tensor nodes, and \(J\) is the number of nodes. With \(m\) denoting the polynomial order and \(L_U\) the constant entering the width rule, define the microcube width
\[
\eta=\min\left\{\frac{1}{2(m+2)},\frac{\lambda_0}{1+L_U}\right\}.
\]
The disjoint reference microcubes are
\[
V_\ell=v_\ell+[-\eta/2,\eta/2]^d,
\qquad \ell=1,\ldots,J.
\]
\end{definitionv}

\begin{definitionv}{synth_3}[Candidate dyadic meshes]
For a sample size \(n\ge1\), dimension \(d\ge1\), and smoothness parameter \(\beta>1\), define the finite candidate mesh collection
\[
\mathcal H_n:=
\left\{2^{-j}:j\in\mathbb N_0,\ 
0\le j\le
\left\lfloor\frac{\log_2 n}{2\beta+d}\right\rfloor
\right\}.
\]
For \(h\in\mathcal H_n\), let \(\mathcal Q_h\) be the dyadic partition of
\(\mathcal X=[0,1]^d\) into cubes of side length \(h\). Each coordinate interval is half-open, with the right endpoint \(1\) included in the last interval of that coordinate. Write \(a_Q\) for the lower-left corner of \(Q\in\mathcal Q_h\).
\end{definitionv}

\begin{algorithmv}{def:mesh-selector}[Count-feasible mesh $\widehat h$]
Given the candidate meshes \(\mathcal H_n\), the observed sample \(\mathcal D_n=((X_i,A_i,Y_i))_{i=1}^n\), the smoothness parameter \(\beta\), and the reference microcubes \(V_\ell\), proceed as follows.
\begin{enumerate}
\item For each \(h\in\mathcal H_n\), form the cube partition \(\mathcal Q_h\). For each partition cube \(Q\), with corner \(a_Q\), define its microcells and treated observation counts by
\[
E_{Q\ell}=a_Q+hV_\ell,
\qquad
N_{Q\ell}=\sum_{i=1}^n\mathbf 1_{\{X_i\in E_{Q\ell},\,A_i=1\}}.
\]
\item Define the set of feasible meshes by
\[
\mathcal F_n
=
\left\{h\in\mathcal H_n:
\min_{Q\in\mathcal Q_h,\,1\le\ell\le J}N_{Q\ell}
\ge h^{-2\beta}\right\},
\]
where \(J\) is the number of reference microcubes.
\item Output the numerically smallest feasible mesh, with zero assigned when the feasible set is empty:
\[
\widehat h=
\begin{cases}
\min\mathcal F_n,&\mathcal F_n\ne\varnothing,\\
0,&\mathcal F_n=\varnothing.
\end{cases}
\]
\end{enumerate}
\end{algorithmv}

\begin{definitionv}{synth_4}[Equal-cell local polynomial coefficients]
Let \(\mathcal D_n=((X_i,A_i,Y_i))_{i=1}^n\) be the observed sample. Fix a feasible mesh \(h\), a partition cube \(Q\in\mathcal Q_h\) with lower-left corner \(a_Q\), the polynomial template vector \(U\), and its \(J\) reference microcubes \(V_\ell\). Define
\[
E_{Q\ell}:=a_Q+hV_\ell,
\qquad
N_{Q\ell}:=\sum_{i=1}^n
\mathbf 1_{\{A_i=1,\ X_i\in E_{Q\ell}\}}.
\]
Feasibility means that \(N_{Q\ell}\ge h^{-2\beta}\) for every partition cube and template microcell, so these counts are positive. The equal-cell empirical Gram matrix is
\[
\widehat G_Q:=
\frac1J\sum_\ell\frac1{N_{Q\ell}}
\sum_{\substack{1\le i\le n\\ A_i=1,\ X_i\in E_{Q\ell}}}
U\!\left(\frac{X_i-a_Q}{h}\right)
U\!\left(\frac{X_i-a_Q}{h}\right)^\top.
\]
Using the polynomial template and its microcell width, this matrix is invertible at every feasible mesh. Define the equal-cell local polynomial coefficient estimate by
\[
\widehat c_Q:=
\widehat G_Q^{-1}
\frac1J\sum_\ell\frac1{N_{Q\ell}}
\sum_{\substack{1\le i\le n\\ A_i=1,\ X_i\in E_{Q\ell}}}
U\!\left(\frac{X_i-a_Q}{h}\right)Y_i.
\]
Thus each microcell receives weight \(1/J\), and the treated observations within that microcell receive equal weights.
\end{definitionv}

\begin{algorithmv}{def:estimator}[Clipped response estimator $\widehat\mu_n$]
Given the selected mesh \(\widehat h\), the polynomial template vector \(U\), and the response bound \(B\), proceed as follows.
\begin{enumerate}
\item If \(\widehat h=0\), output
\[
\widehat\mu_n(x)=0
\qquad(x\in\mathcal X).
\]
\item Otherwise, set
\[
h=\widehat h.
\]
For every \(Q\in\mathcal Q_h\), compute the empirical equal-cell Gram matrix \(\widehat G_Q\) and the equal-cell polynomial coefficient estimate \(\widehat c_Q\).
\item For \(x\in\mathcal X\), let \(Q\) be the unique half-open partition cube containing \(x\), and let \(a_Q\) be its corner. Output
\[
\widehat\mu_n(x)
=
\operatorname{clip}_{[-B,B]}
\left\{
U\left(\frac{x-a_Q}{h}\right)^\top\widehat c_Q
\right\},
\qquad
\operatorname{clip}_{[-B,B]}(t)=\max\{-B,\min\{B,t\}\}.
\]
\end{enumerate}
\end{algorithmv}

\begin{theoremv}{thm:adaptive-minimax}[Adaptive minimax response recovery]
Fix the model parameters subject to the following conditions:
\begin{itemize}
\item \textbf{(Dimension and smoothness.)} \(d\in\mathbb N\), \(d>0\), and \(\beta>1\).
\item \textbf{(Response and smoothness bounds.)} \(B>0\) and \(L>0\).
\item \textbf{(Overlap and density constants.)} \(C\ge1\) and \(0<c_f\le1\).
\item \textbf{(Adaptation interval.)} The real endpoints satisfy \(1<\gamma_{\min}<\gamma_{\max}\), and the adaptation interval is
\[
\Gamma=[\gamma_{\min},\gamma_{\max}].
\]
\end{itemize}
Let \(\mathcal X=[0,1]^d\) be the covariate cube, and let \(\mathcal P_\gamma\) be the observational causal model of \cref{obj:def:model}, with fixed parameters \(d,\beta,B,L,C,c_f,\gamma\), including the propensity-tail envelope in \cref{obj:ass:global-tail}. Define the effective dimension \(D_\gamma\), oracle mesh \(h_{\gamma,n}\), and oracle response-recovery scale \(r_{\gamma,n}\), for \(n\ge1\), by
\[
D_\gamma=\frac{d\gamma}{\gamma-1},
\qquad
h_{\gamma,n}=n^{-1/(2\beta+D_\gamma)},
\qquad
r_{\gamma,n}=h_{\gamma,n}^{\beta}.
\]

The count-selected equal-cell estimator \(\widehat\mu_n\) of \cref{obj:def:estimator} takes \(n,d,\beta,B\) and the observed sample \(\mathcal D_n\) as inputs. There exist \(K>0\) and \(n_0\in\mathbb N\) such that, for every integer \(n\ge\max\{n_0,1\}\), every \(\gamma\in\Gamma\), and every \(P\in\mathcal P_\gamma\),
\[
\mathbb E_{P^{\otimes n}}
\left[
\sup_{x\in\mathcal X}
\left|\widehat\mu_n(x)-\mu_{1,P}(x)\right|
\right]
\le K r_{\gamma,n}.
\]
Here \(P^{\otimes n}\) is the sampling law of \(n\) independent observed triples \((X,A,Y)\), and \(\mu_{1,P}\) is the treated causal response under \(P\).

For every real \(\gamma>1\) and every \(x_0\in\mathcal X\), there exist finite constants \(0<c_\gamma\le K_\gamma\), possibly depending on \(x_0\) and the fixed model parameters, such that, for every integer \(n\ge1\),
\[
\begin{aligned}
c_\gamma r_{\gamma,n}
&\le
\inf_T
\sup_{\substack{P\in\mathcal P_\gamma\\
P\text{ admits a causal completion with}\\
Y(0),Y(1)\in\{0,B\}\ \text{a.s.}}}
\mathbb E_{P^{\otimes n}}
\left|T(\mathcal D_n,x_0)-\mu_{1,P}(x_0)\right|
\\
&\le R_{\mathrm{pt}}(n,\gamma,x_0)
\le R_\infty(n,\gamma)
\le K_\gamma r_{\gamma,n}.
\end{aligned}
\]
The infimum ranges over sample-measurable estimators \(T\); the pointwise and expected supremum minimax risks \(R_{\mathrm{pt}}\) and \(R_\infty\) are those of \cref{obj:def:risks}.
\end{theoremv}

\begin{definitionv}{def:dorn-a3}[Local anti-concentration condition $\mathsf A_3^{\mathrm{Dorn}}$]
The family-level condition
\(\mathsf A_3^{\mathrm{Dorn}}(\mathfrak P,(e_P)_{P\in\mathfrak P})\)
is defined for a family and selected propensity versions satisfying:
\begin{itemize}
\item \textbf{(Family.)} \(\mathfrak P\) is a nonempty family of probability laws with common covariate cube \(\mathcal S\).
\item \textbf{(Versions.)} For each \(P\in\mathfrak P\), \(e_P:\mathcal S\to[0,1]\) is a fixed, pointwise-defined measurable version of \(P(A=1\mid X=\cdot)\).
\item \textbf{(Overlap.)} For each \(P\in\mathfrak P\), \(e_P(X)\in(0,1)\) \(P\)-almost surely.
\end{itemize}
The condition holds when there exist constants \(\rho,\nu>0\) and \(h_0>0\), common to the family, such that, for every \(P\in\mathfrak P\), every \(x_0\in\mathcal S\), and every \(0<h\le h_0\),
\[
P_P\left(
e_P(X)\ge
\rho\sup_{\substack{x\in\mathcal S\\\|x-x_0\|_2\le h}}e_P(x)
\,\middle|\,
\|X-x_0\|_2\le h
\right)>\nu.
\]
Here \(P_P\) denotes probability under \(P\), and the supremum is the pointwise supremum of the selected version. Thus the condition is a property of the family together with its selected propensity versions; its truth can depend on their values on \(P_X\)-null sets. The covariate cube is \(\mathcal S=[-1,1]^d\) in Dorn's source and \(\mathcal S=\mathcal X\) after affine transport. For a single displayed law and propensity, \(\mathsf A_3^{\mathrm{Dorn}}(P)\) abbreviates the condition for the singleton family equipped with that selected version.
\end{definitionv}

\begin{definitionv}{def:thin-strip}[Thin-strip law $P_{\mathrm{strip}}$]
The observed-data law \(P_{\mathrm{strip}}\) is defined by the following construction, under these conditions:
\begin{itemize}
\item \textbf{(Dimension.)} \(d\ge2\).
\item \textbf{(Covariates.)} \(X\) is uniform on \(\mathcal X\).
\item \textbf{(Strip exponent.)} Choose \(a\) satisfying
\[
1<a<1+\frac{d}{\gamma-1}.
\]
\end{itemize}
Define the cube center, radial distance, and its continuous distribution function by
\[
x^\circ=(1/2,\ldots,1/2),
\qquad
R(x)=\|x-x^\circ\|_\infty,
\qquad
F_R(t)=P(R(X)\le t).
\]
Define the strip and selected propensity by
\[
S=\left\{x\in\mathcal X:
|x_2-x_2^\circ|\le |x_1-x_1^\circ|^a\right\},
\qquad
e_\star(x)=
\min\left\{1,F_R(R(x))^{1/(\gamma-1)}
+\frac14\mathbf 1_S(x)\right\}.
\]
Given \(X\), draw
\[
A\sim\operatorname{Bernoulli}(e_\star(X)).
\]
Set
\[
p_0=\min\left\{\frac14,\frac{L}{4B}\right\}.
\]
Draw \(Y(0)\) and \(Y(1)\) conditionally independently given \(X\), each with conditional distribution \(B\operatorname{Bernoulli}(p_0)\), and define
\[
Y=AY(1)+(1-A)Y(0).
\]
Then \(P_{\mathrm{strip}}\) is the law of \((X,A,Y)\).
\end{definitionv}

\begin{propositionv}{prop:strict-enlargement}[Thin-strip enlargement and design degeneration]
Suppose the parameters satisfy:
\begin{itemize}
\item \textbf{(Dimension and smoothness.)} \(d\in\mathbb N\), \(d\ge2\), and \(\beta>1\).
\item \textbf{(Model constants.)} \(B>0\), \(L>0\), \(C\ge1\), and \(0<c_f\le1\).
\item \textbf{(Overlap and strip exponent.)} \(\gamma>1\) and
\[
1<a<1+\frac{d}{\gamma-1}.
\]
\end{itemize}
For every such \(a\), let \(P_{\mathrm{strip}}\) and its displayed selected propensity \(e_\star\) be the construction in \cref{obj:def:thin-strip}, and let \(\mathcal P_\gamma\) denote the model class with the fixed parameters \((d,\beta,B,L,C,c_f,\gamma)\). Then \(P_{\mathrm{strip}}\in\mathcal P_\gamma\), and
\[
P_{\mathrm{strip}}\{e_\star(X)\le t\}\le t^{\gamma-1},
\qquad t\in[0,1].
\]

On the unit cube \(\mathcal X=[0,1]^d\), the numerical anti-concentration clause fails: there are no \(\rho,\nu,h_0>0\) such that, for every \(x_0\in\mathcal X\) and every \(0<h\le h_0\),
\[
\frac{
P_{\mathrm{strip}}\!\left\{
e_\star(X)\ge
\rho\sup_{\substack{x\in\mathcal X\\ \|x-x_0\|_2\le h}}e_\star(x),
\ \|X-x_0\|_2\le h
\right\}
}{
P_{\mathrm{strip}}\{\|X-x_0\|_2\le h\}
}
>\nu.
\]
This failure applies to the numerical clause allowing propensity values equal to one. Consequently, the full singleton-family predicate
\(\mathsf A_{3}^{\mathrm{Dorn}}(P_{\mathrm{strip}})\), evaluated using \(e_\star\), also fails.

Moreover, writing \(x^\circ=(1/2,\ldots,1/2)\), the normalized treated-design second moment in the second coordinate satisfies
\[
\lim_{h\downarrow0}
\frac{
\displaystyle
\int_{\{A=1,\ \|X-x^\circ\|_\infty\le h\}}
\left(\frac{X_2-1/2}{h}\right)^2
\,dP_{\mathrm{strip}}
}{
P_{\mathrm{strip}}\{A=1,\ \|X-x^\circ\|_\infty\le h\}
}
=0.
\]
In particular, \(\mathcal P_\gamma\) contains a law whose full singleton-family Dorn predicate fails for the displayed selected propensity \(e_\star\).
\end{propositionv}

\begin{remarkv}{def:joint-adaptation-handle}[Future work: joint adaptation in $\beta$ and $\gamma$]
A proposed direction for future work is to construct a single response estimator with an adaptive-risk envelope over compact rectangles of the smoothness parameter \(\beta\) and the overlap parameter \(\gamma\). The proposed architecture uses sample splitting to compare equal-cell fits across polynomial degrees and treatment-count-feasible meshes, with the count selector controlling adaptation to overlap and held-out Lepski inequalities controlling adaptation to smoothness. The sample split, degree-and-mesh comparison rule, Lepski threshold, estimator, and envelope functional remain to be specified; estimator existence and the corresponding risk bounds remain open.
\end{remarkv}

\begin{remarkv}{oeq:joint-adaptation}[Joint smoothness and overlap adaptation]
For a compact rectangle of smoothness and overlap parameters \((\beta,\gamma)\) contained in \((1,\infty)^2\), can a procedure adapting along both parameter axes produce a single estimator whose pointwise and expected spatial-supremum risks attain the oracle scale for every pair in the rectangle? What logarithmic penalty, if any, is minimax unavoidable?
\end{remarkv}

\begin{lemmav}{lem:causal-identification}[Identification of the treated response]
Let \(P\in\mathcal P_\gamma\) satisfy the parameter restrictions and causal model conditions of \cref{obj:def:model}, with \(\gamma>1\). Write \(P_X\) for its covariate marginal and \(\mathcal X=[0,1]^d\) for the covariate cube. Then
\[
P_X\{x:e(x)=0\}=0,
\qquad
\mu_1(x)=\mathbb E_P[Y\mid A=1,X=x]
\quad\text{for }P_X\text{-almost every }x.
\]
Moreover, every function \(g:\mathbb R^d\to\mathbb R\) that is continuous on \(\mathcal X\) and satisfies
\[
g(x)=\mathbb E_P[Y\mid A=1,X=x]
\quad\text{for }P_X\text{-almost every }x
\]
obeys
\[
g(x)=\mu_1(x)
\qquad\text{for every }x\in\mathcal X.
\]
\end{lemmav}

\begin{lemmav}{lem:global-holder-extension}[Global Hölder extension bounds]
Fix an integer \(d\ge0\), an integer \(m\ge0\), and an exponent \(0<s\le1\). There exists a constant \(K_{\mathrm{ext}}>0\), depending only on \(d\), \(m\), and \(s\), with the following property. Let \(L\ge0\) and let \(g:[-1,1]^d\to\mathbb R\) satisfy:
\begin{itemize}
\item \textbf{(Intrinsic smoothness.)} The function \(g\) belongs to \(C^m([-1,1]^d)\) in the intrinsic closed-cube sense of \cref{obj:def:holder}, transported to \([-1,1]^d\). Its coordinate derivatives \(D^\alpha g\), for \(|\alpha|\le m\), exist on the interior and admit continuous traces on the closed cube.
\item \textbf{(Derivative bounds.)} With derivatives evaluated through their continuous traces,
\[
\sup_{z\in[-1,1]^d}|D^\alpha g(z)|\le L
\qquad\text{for every }|\alpha|\le m.
\]
\item \textbf{(Hölder continuity.)} For every \(|\alpha|=m\) and all \(z,z'\in[-1,1]^d\),
\[
|D^\alpha g(z)-D^\alpha g(z')|
\le L\|z-z'\|_\infty^{s}.
\]
\end{itemize}
Then there exists \(g_{\mathrm{ext}}\in C^m(\mathbb R^d)\) such that
\[
g_{\mathrm{ext}}(z)=g(z)
\qquad\text{for every }z\in[-1,1]^d,
\]
\[
\sup_{z\in\mathbb R^d}
\|\mathrm D^{k}g_{\mathrm{ext}}(z)\|_{\mathrm{op}}
\le K_{\mathrm{ext}}L
\qquad\text{for every }0\le k\le m,
\]
and
\[
\|\mathrm D^{m}g_{\mathrm{ext}}(z)
-\mathrm D^{m}g_{\mathrm{ext}}(z')\|_{\mathrm{op}}
\le K_{\mathrm{ext}}L\|z-z'\|_\infty^{s}
\qquad\text{for all }z,z'\in\mathbb R^d.
\]
Here \(\mathrm D^{k}\) denotes the Fréchet derivative of order \(k\), with order zero interpreted as the function itself, and \(\|\cdot\|_{\mathrm{op}}\) denotes the operator norm of a \(k\)-linear form on \(\mathbb R^d\) equipped with the supremum norm \(\|\cdot\|_\infty\).
\end{lemmav}

\begin{lemmav}{lem:ordered-mass}[Uniform ordered treated masses]
Fix the model parameters and an overlap interval satisfying
\begin{itemize}
\item \textbf{(Parameter domain.)} \(d\in\mathbb N\), \(d\ge1\), \(\beta>1\), \(B>0\), \(L>0\), \(C\ge1\), and \(0<c_f\le1\).
\item \textbf{(Overlap range.)} \(1<\gamma_{\min}<\gamma_{\max}\), with \(\Gamma=[\gamma_{\min},\gamma_{\max}]\).
\end{itemize}
There exists \(\kappa_\Gamma>0\), depending only on these fixed parameters and interval endpoints, such that the following two bounds hold for every \(\gamma\in\Gamma\) and every causal law \(P\in\mathcal P_\gamma\) satisfying \cref{obj:def:model}.

For every measurable \(E\subseteq\mathcal X=[0,1]^d\), writing \(u=P(X\in E)\),
\[
P(A=1,X\in E)
\ge
\frac{\gamma-1}{\gamma}\,
C^{-1/(\gamma-1)}u^{\gamma/(\gamma-1)}.
\]

For each \(j\in\mathbb N\), set
\[
h=2^{-j},\qquad
m=\lceil\beta\rceil-1,\qquad
D_\gamma=\frac{d\gamma}{\gamma-1}.
\]
Let \(\mathcal Q_h\) be the dyadic cube partition of \(\mathcal X\), with \(|\mathcal Q_h|=(2^j)^d\). Using the degree-\(m\) polynomial template microcells \(E_{Q\ell}\) within each \(Q\in\mathcal Q_h\), define their treated masses and each cube's weakest mass by
\[
q_{Q\ell}=P(A=1,X\in E_{Q\ell}),
\qquad
q_Q=\min_{1\le\ell\le J}q_{Q\ell},
\]
where \(J\) is the number of template microcells. Let \(q_{(k)}\) be the \(k\)th entry of the list of these \(q_Q\)'s sorted in nondecreasing order. Then, for every integer \(1\le k\le(2^j)^d\),
\[
q_{(k)}
\ge
\kappa_\Gamma h^{D_\gamma}k^{1/(\gamma-1)}.
\]
\end{lemmav}

\begin{lemmav}{lem:template-conditioning}[Uniform equal-cell Gram conditioning]
Let \(d,m,n\) be nonnegative integers, let \(\beta\in\mathbb R\), and let \(\mathcal D_n=((X_i,A_i,Y_i))_{i=1}^n\) be any sample with \(X_i\in\mathbb R^d\), \(A_i\in\{0,1\}\), and \(Y_i\in\mathbb R\). Write \(U(x)=(x^\alpha)_{|\alpha|\le m}\) for the vector of total-degree-\(m\) monomials. For the polynomial template's tensor nodes \(v_\ell\) and reference microcubes \(V_\ell\), indexed by \(\ell\in\{0,\ldots,m\}^d\), define
\[
J=(m+1)^d,\qquad
G_0=\frac1J\sum_{\ell\in\{0,\ldots,m\}^d}
U(v_\ell)U(v_\ell)^\top,\qquad
\lambda_0=\inf_{\|b\|_2=1}b^\top G_0b.
\]
Here \(\|b\|_2=(\sum_{|\alpha|\le m}b_\alpha^2)^{1/2}\).

Let \(h=2^{-j}\), with \(j\) a nonnegative integer, be a feasible candidate dyadic mesh: \(h\in\mathcal H_n\), the candidate mesh collection, and
\[
N_{Q\ell}\ge h^{-2\beta}
\quad\text{for every }Q\in\mathcal Q_h
\text{ and }\ell\in\{0,\ldots,m\}^d,
\]
where \(\mathcal Q_h\) is the dyadic cube partition at mesh \(h\). For a cube \(Q\) with corner \(a_Q\), the microcell and its treated count are
\[
E_{Q\ell}=a_Q+hV_\ell,\qquad
N_{Q\ell}=\sum_{i=1}^n
\mathbf1\{A_i=1,\ X_i\in E_{Q\ell}\}.
\]
Define the equal-cell empirical Gram matrix by
\[
\widehat G_Q
=\frac1J\sum_{\ell\in\{0,\ldots,m\}^d}
\frac1{N_{Q\ell}}
\sum_{i=1}^n
\mathbf1\{A_i=1,\ X_i\in E_{Q\ell}\}
U\!\left(\frac{X_i-a_Q}{h}\right)
U\!\left(\frac{X_i-a_Q}{h}\right)^\top.
\]
Then \(\lambda_0>0\), and for every \(Q\in\mathcal Q_h\),
\[
\frac{\lambda_0}{2}\|b\|_2^2
\le b^\top\widehat G_Qb
\quad\text{for every }b=(b_\alpha)_{|\alpha|\le m},
\qquad
\det(\widehat G_Q)\ne0,
\qquad
\|\widehat G_Q^{-1}\|_{\mathrm{op}}
\le\frac2{\lambda_0},
\]
where \(\|M\|_{\mathrm{op}}=\sup_{\|b\|_2=1}\|Mb\|_2\).
\end{lemmav}

\begin{lemmav}{lem:selected-mesh}[Selected mesh and coefficient envelope]
Fix the dimension \(d\), smoothness parameter \(\beta\), model constants \(B,L,C,c_f\), and adaptation endpoints \(\gamma_{\min},\gamma_{\max}\), satisfying
\begin{itemize}
\item \textbf{(Dimension and smoothness.)} \(d\in\mathbb N\), \(d\ge1\), and \(\beta>1\).
\item \textbf{(Model constants.)} \(B>0\), \(L>0\), \(C\ge1\), and \(0<c_f\le1\).
\item \textbf{(Adaptation interval.)} \(1<\gamma_{\min}<\gamma_{\max}\), with
\[
\Gamma=[\gamma_{\min},\gamma_{\max}].
\]
\end{itemize}
There exist constants \(K,K'>0\) and \(n_0\in\mathbb N\) such that the following holds uniformly for every \(n\ge n_0\), \(\gamma\in\Gamma\), and observed law \(P\in\mathcal P_\gamma\), with the displayed fixed model constants. Let the sample have the iid law \(P^{\otimes n}\) specified in \cref{obj:ass:iid}, and use the mesh selector of \cref{obj:def:mesh-selector}. Write
\[
m=\lceil\beta\rceil-1,
\qquad
h_{\gamma,n}=n^{-1/(2\beta+D_\gamma)},
\]
where \(D_\gamma\) is the effective dimension governing the model's overlap-dependent rate, and let \(r_{\gamma,n}\) denote its oracle response-recovery risk scale.

For each partition cube \(Q\) and template microcell \(E_{Q\ell}\), define the population treated mass, its empirical count, and the weakest treated microcell mass by
\[
q_{Q\ell}=P\{A=1,\ X\in E_{Q\ell}\},
\qquad
N_{Q\ell}=\sum_{i=1}^n
\mathbf1\{A_i=1,\ X_i\in E_{Q\ell}\},
\qquad
q_Q=\min_\ell q_{Q\ell}.
\]
There is a measurable sample event \(\mathcal E_n\) such that
\[
P^{\otimes n}(\mathcal E_n^c)\le n^{-2}.
\]
For every sample in \(\mathcal E_n\), the selected mesh satisfies
\[
0<\widehat h\le K h_{\gamma,n},
\qquad
N_{Q\ell}\ge \frac{nq_{Q\ell}}5
\quad
\text{for every }Q\in\mathcal Q_{\widehat h}
\text{ and every }\ell.
\]

Order the selected cubes by increasing \(q_Q\), with ranks starting at \(1\). Conditional on the sampled design \((X_i,A_i)_{i=1}^n\), evaluate all coefficient estimators at the selected mesh of the given sample. For the cube of rank \(k\), every monomial coefficient indexed by a multiindex \(\alpha\) of total degree at most \(m\), and every \(t\in\mathbb R\),
\[
\begin{aligned}
&\mathbb E\!\left[
\exp\!\left\{
t\left(
\widehat c_{Q,\alpha}
-\mathbb E[\widehat c_{Q,\alpha}\mid(X_i,A_i)_{i=1}^n]
\right)
\right\}
\,\middle|\,(X_i,A_i)_{i=1}^n
\right]
\\
&\hspace{2em}\le
\exp\!\left\{
\frac{Kt^2}{2}
\min\!\left(
\widehat h^{2\beta},
n^{-1}\widehat h^{-D_\gamma}k^{-1/(\gamma-1)}
\right)
\right\}.
\end{aligned}
\]
These conditional exponential moments are integrable for every coefficient and every \(t\), and the maximum coefficient-vector error is conditionally integrable. With \(\|\cdot\|_\infty\) denoting the maximum absolute coordinate over those monomial indices, its conditional expectation satisfies
\[
\begin{aligned}
&\mathbb E\!\left[
\max_{Q\in\mathcal Q_{\widehat h}}
\left\|
\widehat c_Q-
\mathbb E[\widehat c_Q\mid(X_i,A_i)_{i=1}^n]
\right\|_\infty
\,\middle|\,(X_i,A_i)_{i=1}^n
\right]
\\
&\hspace{2em}\le
K\widehat h^\beta
\sqrt{1+\max\{0,\log(h_{\gamma,n}/\widehat h)\}}
\le K'r_{\gamma,n}.
\end{aligned}
\]

Moreover, for every exponent \(\alpha>0\), there exists \(K_\alpha>0\) such that the following maximal inequality holds. Let \(\mu\) be a probability law on a measurable space, let \(N\in\mathbb N\), and let \(Z_1,\ldots,Z_N\) be real random variables satisfying
\begin{itemize}
\item \textbf{(Scales.)} \(a>0\) and \(b>0\).
\item \textbf{(Exponential integrability.)} For every \(1\le k\le N\) and \(t\in\mathbb R\), \(\exp(tZ_k)\) is integrable under \(\mu\).
\item \textbf{(Moment envelope.)} For every \(1\le k\le N\) and \(t\in\mathbb R\),
\[
\mathbb E_\mu[\exp(tZ_k)]
\le
\exp\!\left\{
\frac{t^2}{2}\min(a^2,b^2k^{-\alpha})
\right\}.
\]
\end{itemize}
Then
\[
\mathbb E_\mu\!\left[
\max\bigl(\{0\}\cup\{|Z_k|:1\le k\le N\}\bigr)
\right]
\le
K_\alpha a
\sqrt{1+\max\!\left\{0,\log\!\left((b/a)^{2/\alpha}\right)\right\}}.
\]
\end{lemmav}

\begin{definitionv}{def:bounded-lower-pair}[Bounded radial testing pair]
The bounded testing construction is the measure-valued family \(P_{j,h}\) on observations \((X,A,Y)\), indexed by \(j\in\{0,1\}\) and \(n\in\mathbb N\), with parameters \(d\in\mathbb N\), \(\beta,B,L,\gamma\in\mathbb R\), \(x_0\in\mathbb R^d\), and constants \(a,c>0\). Set
\[
\mathcal X:=[0,1]^d,\qquad
h:=c h_{\gamma,n},\qquad
\theta_0:=\min\{1/4,L/(4B)\},
\]
where \(h_{\gamma,n}\) is the oracle mesh for the smoothness and overlap parameters. Define the radial distance, cube-truncated radial function, and common propensity by
\[
R(x):=\|x-x_0\|_\infty,\qquad
F_R(t):=\prod_{i=1}^d
\max\{0,\min\{1,x_{0i}+t\}-\max\{0,x_{0i}-t\}\},
\qquad
e_0(x):=F_R(R(x))^{1/(\gamma-1)}.
\]
For \(x_0\in\mathcal X\) and \(t\ge0\), \(F_R(t)\) is the probability of \(R(X)\le t\) under uniform covariates. Define the smooth product bump and treated success parameters by
\[
\psi(z):=\prod_{i=1}^d
\begin{cases}
\exp\{1-(1-z_i^2)^{-1}\},& |z_i|<1,\\
0,& |z_i|\ge1,
\end{cases}
\qquad
p_{0,h}(x):=\theta_0,\qquad
p_{1,h}(x):=\theta_0+ah^\beta\psi((x-x_0)/h).
\]

The totalized construction uses the positive-part weights
\[
w_0(p):=\max\{1-p,0\},\qquad
w_1(p):=\max\{p,0\}.
\]
For each covariate \(x\), take the product of the three two-point measures with weights
\[
\begin{array}{c|cc}
 & \text{value }0 & \text{value }1\\ \hline
A & w_0(e_0(x)) & w_1(e_0(x))\\
Y(0)/B & w_0(\theta_0) & w_1(\theta_0)\\
Y(1)/B & w_0(p_{j,h}(x)) & w_1(p_{j,h}(x))
\end{array}
\]
where the potential-outcome measures place their masses at \(0\) and \(B\), respectively. Integrate this product measure against Lebesgue measure restricted to \(\mathcal X\), and let \(P_{j,h}\) be its image under
\[
(X,A,Y(0),Y(1))
\longmapsto
\bigl(X,A,AY(1)+(1-A)Y(0)\bigr).
\]

The admissible domain consists exactly of the constants \(a,c>0\) for which
\[
e_0(x),p_{0,h}(x),p_{1,h}(x)\in[0,1]
\]
for every \(n\ge1\) and every \(x\in\mathcal X\), with \(h=c h_{\gamma,n}\). On this domain, \(P_{j,h}\) is the observed law obtained from uniform \(X\) and conditionally mutually independent \(A,Y(0),Y(1)\) given \(X\), with
\[
A\mid X\sim\operatorname{Bernoulli}(e_0(X)),\qquad
Y(0)\mid X\sim B\operatorname{Bernoulli}(\theta_0),\qquad
Y(1)\mid X\sim B\operatorname{Bernoulli}(p_{j,h}(X)),
\qquad
Y=AY(1)+(1-A)Y(0).
\]
For fixed admissible \(a,c\), define the sample-size-indexed collection
\[
\mathcal P_{\mathrm{pair}}
:=\{(P_{0,h},P_{1,h}):n\ge1,\ h=c h_{\gamma,n}\}.
\]
\end{definitionv}

\begin{lemmav}{lem:bounded-lower-pair}[Bounded testing pair guarantees]
Suppose that:
\begin{itemize}
\item \textbf{(Dimension and exponents.)} \(d\in\mathbb N\), \(d>0\), \(\beta>1\), and \(\gamma>1\).
\item \textbf{(Response constants.)} \(B>0\) and \(L>0\).
\item \textbf{(Model constants.)} \(C\ge1\) and \(0<c_f\le1\).
\item \textbf{(Evaluation point.)} \(x_0\in[0,1]^d\).
\end{itemize}
There exist constants \(a,c,c_0>0\), fixed across sample sizes, such that, for every integer \(n\ge1\), with \(h=ch_{\gamma,n}\), where \(h_{\gamma,n}\) is the oracle mesh, the construction in \cref{obj:def:bounded-lower-pair} satisfies
\[
e_0(x),\ p_{0,h}(x),\ p_{1,h}(x)\in[0,1]
\qquad\text{for every }x\in[0,1]^d.
\]
Both observed laws \(P_{0,h}\) and \(P_{1,h}\) belong to \(\mathcal P_\gamma\), the model class with parameters \(d,\beta,B,L,C,c_f,\gamma\). Under each of their constructed causal completions,
\[
Y(0),Y(1)\in\{0,B\}
\qquad\text{almost surely}.
\]
Moreover, \(P_{0,h}\ll P_{1,h}\), the log likelihood ratio
\(\log(dP_{0,h}/dP_{1,h})\) is integrable under \(P_{0,h}\), and
\[
n\,\operatorname{KL}(P_{0,h},P_{1,h})\le\frac18.
\]
For the chosen Hölder response representatives \(\mu_{1,P_{0,h}}\) and \(\mu_{1,P_{1,h}}\) associated with their model-class memberships,
\[
\left|\mu_{1,P_{1,h}}(x_0)-\mu_{1,P_{0,h}}(x_0)\right|
\ge c_0 r_{\gamma,n},
\qquad
r_{\gamma,n}=h_{\gamma,n}^{\beta},
\]
where \(r_{\gamma,n}\) is the oracle response-recovery risk scale.
\end{lemmav}

\begin{lemmav}{lem:fixed-cube-holder-completion}[Hölder completion under affine scaling]
Fix an integer \(d\ge 1\) and a real number \(\beta>1\), and set \(m=\lceil\beta\rceil-1\). There exists a constant \(K_{d,\beta}>0\), depending only on \(d\) and \(\beta\), such that the following holds for every \(M_0,L_0\in\mathbb R\) and \(u:\mathbb R^d\to\mathbb R\) satisfying:
\begin{itemize}
\item \textbf{(Nonnegative bounds.)} \(M_0\ge0\) and \(L_0\ge0\).
\item \textbf{(Intrinsic regularity.)} \(u\in C^m([-1,1]^d)\), with interior derivatives through order \(m\) admitting continuous traces on the closed cube.
\item \textbf{(Top derivative modulus.)} For every multi-index \(\alpha\) with \(|\alpha|=m\) and every \(x,y\in[-1,1]^d\), the coordinate derivatives, interpreted through their continuous boundary traces, satisfy
\[
|D^\alpha u(x)-D^\alpha u(y)|
\le L_0\left(\sum_{i=1}^d(x_i-y_i)^2\right)^{(\beta-m)/2}.
\]
\item \textbf{(Response bound.)} \(|u(x)|\le M_0\) for every \(x\in[-1,1]^d\).
\end{itemize}
For the affine transport
\[
T(x)=2x-\mathbf1,\qquad x\in[0,1]^d,
\]
the transported function satisfies
\[
u\circ T\in\mathcal H^\beta\!\left(K_{d,\beta}(M_0+L_0)\right)
\]
on \([0,1]^d\), with the intrinsic Hölder ball defined in \cref{obj:def:holder}.
\end{lemmav}

\begin{propositionv}{prop:dorn-a3-free-corollary}[Transfer of Dorn source families]*
Fix parameters satisfying:
\begin{itemize}
\item \textbf{(Dimension and smoothness.)} \(d\in\mathbb N\), \(d\ge1\), and \(\beta>1\).
\item \textbf{(Source constants.)} \(q>3\) and \(M,\sigma_{\min},C,L_0>0\).
\item \textbf{(Adaptation range.)} \(1<\gamma_{\min}<\gamma_{\max}\), with
\[
\Gamma=[\gamma_{\min},\gamma_{\max}].
\]
\end{itemize}

For each \(\gamma\in\Gamma\), let \(\mathcal D_\gamma\) be the maximal envelope of source-law tuples consisting of an observed law \(P\), treated and control mean functions, and a selected propensity \(e_P\), satisfying the following retained restrictions from \citet[Section 2, Assumptions 1--2, pp.~3--7]{Dorn2026GlobalOverlap}: \(X\) is uniform on \([-1,1]^d\); the selected propensity belongs to \((0,1)\) almost surely; conditional on \(X,A\), the outcome is Gaussian with variance \(\sigma_{\min}^2\); both conditional mean functions belong to \(\Sigma^{\mathrm D}(\beta,L_0)\); and
\[
\mathbb E_P\!\left[|Y|^q\mid X,A=1\right]\le M^q,
\qquad
\operatorname{Var}_P(\mu_P(X))\le M,
\qquad
\operatorname{Var}_P(Y\mid X,A)\ge\sigma_{\min}^2,
\]
\[
P\{e_P(X)\le t\}\le Ct^{\gamma-1}
\qquad (0\le t\le1).
\]
Here \(\mu_P\) is the selected treated mean. The source class \(\Sigma^{\mathrm D}(\beta,L_0)\) consists of functions in \(C^m([-1,1]^d)\), with continuous boundary traces of their derivatives, whose derivatives of order
\[
m=\max\{k\in\mathbb Z:k<\beta\}
\]
have \((\beta-m)\)-Hölder seminorm at most \(L_0\), using the source's convention.

Let \(T\), the affine transport from the unit cube to the source cube, and the common response bound \(B_*\) be
\[
T(x)=2x-\mathbf1,
\qquad
B_*=\max\{M,\sigma_{\min}\}.
\]
Define the common transported Hölder radius and the interpolation constant by
\[
L_*=
\sup_{\substack{u\in\Sigma^{\mathrm D}(\beta,L_0)\\
                 \sup_{x\in[-1,1]^d}|u(x)|\le M}}
\|u\circ T\|_{\mathcal H^\beta},
\]
\[
C_{d,\beta}=
\sup_{\substack{M'>0,\ L_0'>0\\
                 u\in\Sigma^{\mathrm D}(\beta,L_0')\\
                 \sup_{x\in[-1,1]^d}|u(x)|\le M'}}
\frac{\|u\circ T\|_{\mathcal H^\beta}}{M'+L_0'}.
\]
The functional \(\|\cdot\|_{\mathcal H^\beta}\) is the full Hölder-radius functional on \([0,1]^d\), including the function, lower derivatives, and top-derivative modulus. Then
\[
0<C_{d,\beta}<\infty,
\qquad
0<L_*\le C_{d,\beta}(M+L_0)<\infty.
\]

Every member of \(\mathcal D_\gamma\), after replacing \(X\) by \(T^{-1}(X)\), belongs to the global-tail causal class \(\mathcal P_\gamma\) with constants
\[
B=B_*,
\qquad L=L_*,
\qquad C=C,
\qquad c_f=1.
\]
In particular, its transported observed law admits a consistency-and-exchangeability causal completion with these constants.

Let \(\widehat\mu_n\) be the equal-cell estimator of \cref{obj:def:estimator}, with smoothness \(\beta\) and response bound \(B_*\), applied to \(n\) independent transported observations. Write \(r_{\gamma,n}\), the oracle risk scale, as
\[
r_{\gamma,n}=n^{-\beta/(2\beta+d\gamma/(\gamma-1))}
\qquad(n\ge1).
\]
There exist \(K>0\) and \(n_0\in\mathbb N\) such that, for every \(n\ge n_0\), every \(\gamma\in\Gamma\), and every source tuple in \(\mathcal D_\gamma\),
\[
\mathbb E_{P^{\otimes n}}
\left[
\sup_{x\in[0,1]^d}
\left|\widehat\mu_n(x)-\mu_P(T(x))\right|
\right]
\le K r_{\gamma,n}.
\]
Every source family \(\mathfrak P\) satisfying the cited Assumptions 1--2 with the fixed displayed constants is nonempty and contained in \(\mathcal D_\gamma\). The same \(K,n_0\) therefore satisfy
\[
\sup_{P\in\mathfrak P}
\mathbb E_{P^{\otimes n}}
\left[
\sup_{x\in[0,1]^d}
\left|\widehat\mu_n(x)-\mu_P(T(x))\right|
\right]
\le K r_{\gamma,n}
\qquad(n\ge n_0).
\]

For every \(\varepsilon>0\), there exist \(c(\varepsilon)>0\) and \(n_0\in\mathbb N\) such that, simultaneously for all \(n\ge n_0\), \(\gamma\in\Gamma\), and source tuples in \(\mathcal D_\gamma\),
\[
P^{\otimes n}\!\left\{
\left\|\widehat\mu_n\circ T^{-1}-\mu_P\right\|_{L^\infty(P)}
>c(\varepsilon)r_{\gamma,n}
\right\}\le\varepsilon.
\]
Here the norm is the essential supremum under the source covariate law, and the estimator is computed from the transported sample.

If \(d\ge2\), then for each \(\gamma\in\Gamma\) there exist positive finite constants \(C',M',\sigma_{\min}',L_0'\) and a source tuple with uniform covariates, Gaussian conditional outcomes, and a selected pointwise propensity representative satisfying all the retained restrictions with these constants, for which
\[
\neg\,
\mathsf A_{3}^{\mathrm{Dorn}}
\bigl(\{P\},(e_P)\bigr),
\]
where the family-level condition is defined in \cref{obj:def:dorn-a3}. The constants in this existence assertion are selected separately for each \(\gamma\).

For any nonempty family \(\mathfrak P\) of source-law tuples sharing a common treated regression curve \(\mu^\circ\), meaning
\[
\mu_P(x)=\mu^\circ(x)
\qquad(P\in\mathfrak P,\ x\in\mathbb R^d),
\]
the minimax pointwise miss-probability value at any evaluation point \(x_0\), sample size \(n\in\mathbb N\), and positive threshold \(cr\), with \(c,r>0\), equals zero.

Finally, for every \(\gamma>1\) and every \(x_0\in[0,1]^d\), there exists \(c_\gamma>0\) such that
\[
c_\gamma r_{\gamma,n}
\le R_{\mathrm{pt}}(n,\gamma,x_0)
\qquad(n\ge1),
\]
where \(R_{\mathrm{pt}}\) is the pointwise minimax expected absolute-error risk of \cref{obj:def:risks} over the causal class with constants \(B_*,L_*,\max\{C,1\}\), and \(c_f=1\). The displayed lower bound is first established for the two-valued potential-outcome subclass risk; monotonicity under class inclusion transfers it to \(R_{\mathrm{pt}}\). The constant \(c_\gamma\) may depend on the evaluation point and the fixed parameters.
\end{propositionv}
% CAUSALSMITH-CITED-SCOPE-BEGIN prop:dorn-a3-free-corollary
\verificationfootnotetext{\textbf{Formalization scope.} This result uses the cited conclusion \citep{Dorn2026GlobalOverlap} (Section 2 Setting and Notation, Assumptions 1--2; author PDF pages 3--7). The cited source proof is not formalized here: Lean verifies its use and all remaining steps. If that cited conclusion is false, the portions of this result that depend on it are not certified.}
% CAUSALSMITH-CITED-SCOPE-END prop:dorn-a3-free-corollary

\begin{lemmav}{lem:arbitrary-family-lower-obstruction}[Known regression minimax obstruction]
Fix parameters satisfying
\begin{itemize}
\item \textbf{(Dimension.)} \(d\in\mathbb N\) and \(d\ge1\).
\item \textbf{(Smoothness.)} \(\beta>0\) and \(L_0>0\).
\item \textbf{(Overlap.)} \(\gamma\in\mathbb R\) and \(\gamma>1\).
\end{itemize}
Let \(\mathfrak P\) be any nonempty source family equipped with selected treated and control regression functions and propensity representatives. Suppose its selected treated regression \(\mu_P\) is the same known function \(\mu^\circ:\mathbb R^d\to\mathbb R\) throughout the family:
\[
\mu_P(x)=\mu^\circ(x)
\qquad
\text{for every }P\in\mathfrak P
\text{ and every }x\in\mathbb R^d.
\]
For a sample size \(n\in\mathbb N\), let \(\mathcal D_n\) denote the observed sample and \(P^{\otimes n}\) its product sampling law. Then, for every evaluation point \(x_0\in\mathbb R^d\), every \(n\in\mathbb N\), and every \(c,r_n>0\), where \(r_n\) is an arbitrary positive comparison scale,
\[
\inf_{T\ \mathrm{measurable}}
\sup_{P\in\mathfrak P}
P^{\otimes n}
\left\{\left|T(\mathcal D_n)-\mu_P(x_0)\right|>c r_n\right\}
=0.
\]
The infimum ranges over all measurable real-valued sample estimators \(T\), and the minimax value is taken in the extended nonnegative reals.

In particular, consider the singleton family whose observed law has mutually independent coordinates
\[
X\sim\operatorname{Unif}([-1,1]^d),
\qquad
A\sim\operatorname{Bernoulli}(1/2),
\qquad
Y\sim N(0,1),
\]
with selected treated and control regression functions identically zero and selected propensity \(e_P(x)=1/2\). There exist common constants \(q>3\), \(M>0\), and \(C>0\) for which this law satisfies all of the following:
\begin{itemize}
\item \textbf{(Selected propensity.)} The selected propensity is a measurable version of the conditional treatment probability and lies in \((0,1)\) almost surely.
\item \textbf{(Moment and variance bounds.)} The treated conditional absolute \(q\)-moment is integrable and satisfies
\[
\mathbb E_P\!\left[|Y|^q\mid X=x,A=1\right]\le M^q
\quad\text{for }P_X\text{-almost every }x,
\qquad
\mathbb E_P\!\left[
\left(\mu_P(X)-\mathbb E_P[\mu_P(X)]\right)^2
\right]\le M.
\]
\item \textbf{(Source smoothness.)} Both selected regression functions belong to the source top-derivative seminorm class \(\Sigma^{\mathrm D}(\beta,L_0)\) specified in \cref{obj:prop:dorn-a3-free-corollary}.
\item \textbf{(Global overlap.)} The selected propensity satisfies
\[
P_X\{e_P(X)\le t\}\le C t^{\gamma-1}
\qquad\text{for every }t\in[0,1].
\]
\item \textbf{(Conditional Gaussian outcomes.)} For \(P_X\)-almost every \(x\), the conditional outcome distribution in each treatment arm is Gaussian with its selected regression mean and variance one.
\item \textbf{(Common local anti-concentration.)} The singleton satisfies
\[
\mathsf A_{3}^{\mathrm{Dorn}}
\bigl(\mathfrak P,(e_P)_{P\in\mathfrak P}\bigr),
\]
with the local ball-ratio inequalities and common positive constants specified in \cref{obj:def:dorn-a3}.
\end{itemize}
At the cube center \(x_0=0\), for every \(n\in\mathbb N\) and every \(c,r_n>0\), this singleton's displayed minimax miss probability equals zero.
\end{lemmav}

\begin{propositionv}{lem:dorn-rate-scope}[Fixed-family regression rate scope]*
Fix an integer \(d\ge1\), constants \(\beta,L_0,M,\sigma,C>0\), and \(q>3\). For every fixed finite \(\gamma_0>1\), consider a nonempty source family \(\mathfrak P\) of observed-data laws, equipped with selected treated and control regression curves and selected propensity representatives \((e_P)_{P\in\mathfrak P}\). Suppose that:
\begin{itemize}
\item \textbf{(Design and propensity.)} Under every \(P\in\mathfrak P\), \(X\) is uniform on \([-1,1]^d\), and the selected measurable representative satisfies
\[
e_P(X)=P(A=1\mid X)\in(0,1)
\qquad P\text{-almost surely}.
\]
\item \textbf{(Moment and variance.)} Writing \(\mu_P\) for the selected treated regression curve, every member satisfies
\[
\mathbb E_P\!\left[|Y|^q\mid X,A=1\right]\le M^q,\qquad
\operatorname{Var}_P\!\left(\mu_P(X)\right)\le M,\qquad
\operatorname{Var}_P(Y\mid X,A)\ge\sigma^2,
\]
with the conditional inequalities holding almost surely in each applicable arm.
\item \textbf{(Gaussian regression and smoothness.)} Almost surely with respect to the covariate marginal, the conditional outcome distributions in the treated and control arms are Gaussian with their respective selected regression means and common variance \(\sigma^2\). Both selected regression curves belong to the source seminorm class \(\Sigma^{\mathrm D}(\beta,L_0)\): their derivatives of order
\[
m=\max\{k\in\mathbb Z:k<\beta\}
\]
satisfy Dorn's Hölder-seminorm bound \(L_0\), with exponent \(\beta-m\).
\item \textbf{(Global overlap tail.)} For every \(P\in\mathfrak P\) and \(t\in[0,1]\),
\[
P\{e_P(X)\le t\}\le Ct^{\gamma_0-1}.
\]
\item \textbf{(Family-level local condition.)} On the original cube \([-1,1]^d\), the selected propensity representatives satisfy
\[
\mathsf A_{3}^{\mathrm{Dorn}}
\bigl(\mathfrak P,(e_P)_{P\in\mathfrak P}\bigr)
\]
as defined in \cref{obj:def:dorn-a3}, including its common positive constants \(\rho,\nu,h_0\) and uniform local-ball mass-ratio inequality.
\end{itemize}

Define the fixed-setup comparison scale by
\[
r^*_{\mu,n}
:=n^{-\beta/(2\beta+d\gamma_0/(\gamma_0-1))}.
\]
There exists a regression estimator \(\widehat\mu_n(\mathcal D_n,x)\), where \(\mathcal D_n\) is the sample of \(n\) independent observations, selected after the family, its common constants, and its propensity representatives have been fixed, such that its dependence on the sample is measurable for every \(n\) and \(x\). For every sample size \(n\), member \(P\in\mathfrak P\), and real threshold \(R\), the event
\[
\left\{
\|\widehat\mu_n-\mu_P\|_{L^\infty(P)}
>\max\{Rr^*_{\mu,n},0\}
\right\}
\]
is measurable, where the norm is the essential supremum with respect to the covariate marginal of \(P\). For every \(\varepsilon>0\), there exists a finite constant \(c(\varepsilon)>0\) such that
\[
\limsup_{n\to\infty}\sup_{P\in\mathfrak P}
P^{\otimes n}\!\left\{
\|\widehat\mu_n-\mu_P\|_{L^\infty(P)}
>c(\varepsilon)r^*_{\mu,n}
\right\}
\le\varepsilon,
\]
where \(P^{\otimes n}\) is the product sampling law. The estimator and \(c(\varepsilon)\) may depend on the fixed setup. This upper comparison is the assertion of \citet[Section 2, Assumptions 1--3, Theorem 1(ii)]{Dorn2026GlobalOverlap}.

For every \(\gamma_0>1\), the following two pointwise conclusions also hold. First, let \(\mathfrak P\) be any nonempty source family whose selected treated regression curves equal a common known function \(\mu^\circ\) at every point:
\[
\mu_P(x)=\mu^\circ(x)
\qquad(P\in\mathfrak P,\ x\in\mathbb R^d).
\]
For every evaluation point \(x_0\in\mathbb R^d\), sample size \(n\), and constants \(c,r>0\), its extended nonnegative minimax miss risk is
\[
\inf_{T\ \mathrm{sample\text{-}measurable}}
\sup_{P\in\mathfrak P}
P^{\otimes n}\!\left\{
|T(\mathcal D_n)-\mu_P(x_0)|>cr
\right\}
=0.
\]
Here the infimum ranges over all measurable real-valued functions of the sample.

Second, there exist constants \(q'>3\) and \(M',C'>0\) for which the singleton Gaussian experiment
\[
X\sim\operatorname{Unif}([-1,1]^d),\qquad
A\sim\operatorname{Bernoulli}(1/2),\qquad
Y\sim N(0,1),
\]
with \(X,A,Y\) mutually independent and selected treated mean, control mean, and propensity respectively \(0,0,1/2\), satisfies the displayed design, propensity, moment, variance, Gaussian, smoothness, and global-tail conditions with parameter tuple
\[
(\beta,q',M',1,C',L_0,\gamma_0),
\]
and satisfies the family-level condition of \cref{obj:def:dorn-a3} on the original cube. For this singleton, at \(x_0=0\), every sample size \(n\) and every \(c,r>0\) satisfy
\[
\inf_{T\ \mathrm{sample\text{-}measurable}}
P^{\otimes n}\!\left\{|T(\mathcal D_n)|>cr\right\}
=0.
\]
\end{propositionv}
% CAUSALSMITH-CITED-SCOPE-BEGIN lem:dorn-rate-scope
\verificationfootnotetext{\textbf{Formalization scope.} This result uses the cited conclusion \citep{Dorn2026GlobalOverlap} (Section 2 Notation and Setting, Assumptions 1--3, and Theorem 1 (author PDF pp. 3--7 and 12); Lemma 11 (p. 23); Lemma 13 (p. 24); proof of Theorem 1 (pp. 27--28); Section 5 Conclusion (p. 28)). The cited source proof is not formalized here: Lean verifies its use and all remaining steps. If that cited conclusion is false, the portions of this result that depend on it are not certified.}
% CAUSALSMITH-CITED-SCOPE-END lem:dorn-rate-scope
